\documentclass[12pt]{article}
\newcommand{\ul}[1]{\underline{#1}}
\usepackage{amsmath}
\usepackage{graphicx}
\usepackage{float}

\setlength{\textwidth}{6.5in}
\setlength{\textheight}{9.0in}
\setlength{\oddsidemargin}{-0.25in}
\setlength{\evensidemargin}{-0.25in}
\setlength{\topmargin}{-0.75in}

\begin{document}
\noindent{\large\bf EGCP-3210 \hfill Computer Architecture \hfill Spring 2026} \\
Homework 6 \hfill {\bf Brandon Aikman} \\

% Homework 6
\begin{enumerate}
	% 1.
	\item A two-way set associative cache has lines of 16 bytes and a total data size of 8 kbytes. The 64-Mbyte main memory is byte addressable. Show the format of main memory address in terms of tag, index and word. (10 pts)
	
	% 2.
	\item For the hexadecimal main memory address 111111, 666666, BBBBBB, show the following information, in hexadecimal format, assume there are 16K blocks in the cache, each block has 4 words. (5 pts each; 15 pts total)
    \begin{enumerate}
        \item Tag , index, and word value for a direct-mapped cache.
        \item Tag and word value for a fully associative cache.
        \item Tag, set and word value for a two-way set associative.
    \end{enumerate}

    % 3.
    \item Consider a machine with a byte addressable main memory of 64k bytes and block size of 8 bytes. Assume that a direct mapped cache consisting of 32 lines is used with this machine. (5 pts each; 15 pts total)
    \begin{enumerate}
        \item Into what line would bytes with each of the following addresses be stored? \\
            0001 0001 0001 1011 \\
            1100 0011 0011 0100 \\
            1101 0000 0001 1101 \\
            1010 1010 1010 1010 \\
        \item Suppose the byte with address 0001 1010 0001 1010 is stored in the cache. What are the addresses of the other bytes stored along with it?
        \item How many total bits of this cache. (including a validation bit for each cache line)
    \end{enumerate}

    % 4.
    \item Consider a cache of 4 lines of 16 bytes each. Main memory is divided into blocks of 16 bytes each. That is, block 0 has bytes with address 0 through 15, and so on. The Least Recently Used replacement scheme is used. Now consider a program that access memory in the following sequence of addresses:\\
    Once: 63 through 70 \\
    Loop ten times: 15 through 32; 80 through 95 (20 pts)
    \begin{enumerate}
        \item Suppose the cache is organized as direct mapped. Compute the hit ratio.
        \item Suppose the cache is organized as two-way set associative. Compute the hit ratio.
    \end{enumerate}
    
\end{enumerate}

\end{document}